\subsection{Activity patterns}
\label{sec:r-activity}

% Owner: C
\paragraph{Burstiness and memory.}
\expect{[Expected] Inter-event times of both PS and CM are heavy-tailed and have
positive burstiness $B$ and memory $M$. Under P1 the heavy tail is largely
removed, so $B_{\mathrm{real}}-B_{\mathrm{null}}>0$. Contact bursts in PS are
bounded by the school timetable, while CM shows longer waiting times (hours to
days) and a stronger tail.}

\begin{figure}[t]
\centering
\figph{\columnwidth}{3.6cm}{inter-event time distributions, real vs.\ P1\\(PS, CM)}
\caption{Inter-event time distributions on log--log axes. (Placeholder.)}
\label{fig:iet}
\end{figure}

\begin{table}[t]
\centering
\small
\begin{tabular}{@{}lcccc@{}}
\toprule
 & \multicolumn{2}{c}{PS} & \multicolumn{2}{c}{CM} \\
 & real & P1 & real & P1 \\
\midrule
Burstiness $B$ & $>0$ & $\approx$ lower & $>0$ & $\approx$ lower \\
Memory $M$ & $>0$ & $\approx 0$ & $>0$ & $\approx 0$ \\
\bottomrule
\end{tabular}
\caption{Expected qualitative outcome of M2a (values to be measured; mean $\pm$
std over $\geq 20$ P1 samples).}
\label{tab:burst}
\end{table}

% Owner: C
\paragraph{Persistence and turnover.}
\expect{[Expected] Many contacts in PS are repeated within the same day (within
class), so the edge repetition rate is high and consecutive-snapshot Jaccard is
moderate, with drops at class changes. In CM most pairs interact only a few
times, so tie lifetimes are short and turnover is high. P1 keeps the static graph
and therefore the repetition rate, but changes lifetimes and Jaccard.}

\paragraph{Verdict on H2 and H3.}
\outline{Supported / not supported / depends on network type or window.}
\expect{[Expected] H2 supported in both. H3 supported in PS, only partly in CM.}
